\subsection{Hermitian endomorphisms}

DEFINITION 5. --- Let E be a hilbertian space and let \(u \in \mathcal{L}(\mathrm{E})\) . We say that \(u\) is hermitian if \(u^* = u\) .

Let \(\mathcal{H}(\mathrm{E})\) denote the set of all hermitian elements of \(\mathcal{L}(\mathrm{E})\) ; this is a vector subspace of the vector space \(\mathcal{L}(\mathrm{E})_{[\mathbb{R}]}\) over \(\mathbf{R}\) which is deduced from \(\mathcal{L}(\mathrm{E})\) by restricting the scalars.

To each \(u \in \mathcal{L}(\mathrm{E})\) , we associated (V, p.~16, cor. 2) a sesquilinear form \(\Phi_u : (x, y) \mapsto \langle x | u(y) \rangle\) on \(\mathbf{E} \times \mathbf{E}\) . We have

\[
\Phi_ {u ^ {*}} (x, y) = \overline {{{{\Phi_ {u} (y , x)}}}} \quad (x, y \text {in E});\tag{14}
\]

consequently, \(u\) is hermitian if and only if the form \(\Phi_u\) is hermitian. When \(K\) is \(C\) , it is enough to assume that \(\Phi_u(x, x) = \langle x | u(x) \rangle\) is real for all \(x \in E\) (V, p.~2, Remark).

Let \(u \in \mathcal{L}(\mathrm{E})\) . We have seen (V, p.~16, cor. 2) that the norm of \(u\) can be calculated by the formula

\[
\| u \| = \sup _ {\| x \| \leqslant 1, \| y \| \leqslant 1} \left| \Phi_ {u} (x, y) \right|.\tag{15}
\]

When u is hermitian, we have the following result :

PROPOSITION 9. --- For every hermitian endomorphism u of E, we have

\[
\| u \| = \sup _ {\| x \| \leqslant 1} | \langle x | u (x) \rangle |.\tag{16}
\]

Put \(\Phi = \Phi_{u}\) and \(c = \sup_{\| x\| \leqslant 1}|\Phi (x,x)|\) , then evidently \(c\leqslant \| u\|\) . Let \(x,y\) be in E such that \(\| x\| \leqslant 1\) , \(\| y\| \leqslant 1\) . Then

\[
\Phi (x + y, x + y) = \Phi (x, x) + \Phi (y, y) + 2 \mathscr {R} \Phi (x, y),
\]

hence

\[
4 \mathscr {R} \Phi (x, y) = \Phi (x + y, x + y) - \Phi (x - y, \dot {x} - y);
\]

but \(|\Phi(t, t)| \leqslant c \|t\|^2\) for all \(t \in \mathbf{E}\) , thus

\[
4 \left| \mathcal {R} \Phi (x, y) \right| \leqslant c \left(\| x + y \| ^ {2} + \| x - y \| ^ {2}\right) = 2 c \left(\| x \| ^ {2} + \| y \| ^ {2}\right) \leqslant 4 c.
\]

Let \(a = \Phi(x, y)\) ; there exists a complex number \(\lambda\) with absolute value 1 such that \(\lambda a = |a|\) . Replacing \(y\) by \(\lambda y\) in the preceding inequality, we get \(|\Phi(x,y)| \leqslant c\) . By (15), \(\|u\| \leqslant c\) and the proposition follows. Q.E.D.

Evidently every hermitian endomorphism is normal. Conversely :

PROPOSITION 10. --- Suppose K is C. Let \(u \in \mathcal{L}(\mathrm{E})\) . Then there exists a unique pair \((h_1, h_2)\) of hermitian endomorphisms of E, such that \(u = h_1 + ih_2\) . In order that u is normal, it is necessary and sufficient that \(h_1\) and \(h_2\) commute.

For, the relation \(\ll u = h_1 + ih_2\) , \(h_1^* = h_1\) , \(h_2^* = h_2\) is equivalent to

\[
\ll h _ {1} = \frac {1}{2} (u + u ^ {*}) \text { and } h _ {2} = \frac {i}{2} (u ^ {*} - u) \gg .
\]

In addition, we have \(h_1h_2 - h_2h_1 = \frac{i}{2} (uu^* - u^*u)\) . This proves prop. 10.

PROPOSITION 11. --- Let \(p \in \mathcal{L}(\mathbf{E})\) . In order that \(p\) is the orthoprojector from \(\mathbf{E}\) onto a closed vector subspace of \(\mathbf{E}\) , it is necessary and sufficient that \(p^2 = p = p^*\) .

Suppose \(p^{2} = p\) . Let M be the image of p and N its kernel. E is the topological direct sum of M and N. In order that p is an orthoprojector, it is necessary and sufficient that M is orthogonal to N, that is to say that we have \(\langle p(x)|y - p(y)\rangle = 0\) for all x, y in E. This latter relation is equivalent to \(p = p^{*}p\) , and implies that \(p^{*} = (p^{*}p)^{*} = p^{*}p = p\) ; conversely if \(p^{*} = p\) , we have \(p = p^{2} = p^{*}p\) .

